\begin{claim}{F.14}{Proved}{An injective seven-point comparison and its exact remaining defect}
Retain the based map $\phi$ and the involution $R$ of F.12.
Label its eight target points
\[
             0,\ \infty,\ v_0=1,\ v_1,\ldots,v_5 .
                                                               \tag{F.14.1}
\]
Let $q_{v_0}$ forget the ordinary point $v_0$, retaining
both branch labels $0,\infty$.  The actual composite
\[
 \widehat\psi=(q_{v_0})_*\widehat\phi:
       \widehat{\pi_1(Y,s_0)}
                     \longrightarrow\widehat{P\Gamma}_{0,7}
                                                               \tag{F.14.2}
\]
is injective.  With $v_1$ renormalized to one, its four
remaining coordinates and those of the second composite
$\psi_R=q_{v_0}R\phi$ are exactly
\[
 \begin{split}
 \psi(s)&=(v_2/v_1,\ v_3/v_1,\ v_4/v_1,\ v_5/v_1),\\
 \psi_R(s)&=(v_3,\ v_3/v_1,\ v_3/v_5,\ v_3/v_4).
 \end{split}                                                   \tag{F.14.3}
\]
Both maps have the same specified target basepoint.
Their second coordinates agree as actual functions.

Forgetting either branch point $0$ or $\infty$ from
(F.14.1) instead kills each of the two explicitly
whiskered source meridians at $s=0,\infty$ below.
Each meridian generates a full closed $\widehat\Z$
in the source.  Both branch-forgetting composites
factor through the specified filling
\[
 \widehat{\pi_1(Y)}
       \longrightarrow
 \widehat{\pi_1\bigl(\mathbf P^1\setminus\{\pm\mu_5\}\bigr)} .
                                                               \tag{F.14.4}
\]

There is also an exact description of the information
retained by (F.14.3).  Put $P_8=\widehat{P\Gamma}_{0,8}$,
write $K_i$ for its actual kernel of forgetting $v_i$,
and define
\[
 \begin{split}
 \mathcal E_7&=\{u:\widehat\psi(u)=\widehat\psi_R(u)\},\\
 B&=K_0\cap K_2,\\
 d(u)&=\widehat\phi(u)^{-1}R_*\widehat\phi(u)
                         \qquad(u\in\mathcal E_7).
 \end{split}                                                   \tag{F.14.5}
\]
Then $d$ takes values in $B$, satisfies
\[
 \begin{split}
 R_*d(u)&=d(u)^{-1},\\
 d(uv)&=\widehat\phi(v)^{-1}d(u)\widehat\phi(v)d(v),\\
 \mathcal E&=\{u\in\mathcal E_7:d(u)=1\},
 \end{split}                                                   \tag{F.14.6}
\]
where $\mathcal E$ is the full residual equalizer F.12.5.
In the actual point-pushing identification $K_0=\widehat F_6$,
$B$ is precisely the closed normal subgroup generated by
the meridian about $v_2$.  These statements retain both
distinct forgetting composites; they do not cancel
$q_{v_0}$ from the equalizer.
\end{claim}

\begin{proof}
\emph{Filling the two order-five poles.}
Let $S_{12}$ be the sphere punctured at the labelled
configuration
\[
                  0,\infty,\mu_5,t_0\mu_5,
\]
and let $S_{10}$ be obtained by filling the first two
punctures.  The named rotation $d(z)=\zeta z$ acts on
both surfaces.  The two fixed poles have become
unmarked ramification points on $S_{10}$.
After removal of those ramification points, both
covering problems are literally the same unramified
cover
\[
 S_{12}\longrightarrow
       X=\mathbf P^1\setminus\{0,\infty,1,t_0^5\},
                           \qquad z\longmapsto z^5 .
                                                               \tag{F.14.7}
\]
Its four meridians have the same monodromy values
$1,-1,0,0$ in $C_5$ as F.11.8.

Write $N_{12}$ and $N_{10}$ for the discrete full
normalizers of this specified $C_5$ in the two
mapping class groups.  The actual covering normalizer
sequences of F.7.5 form the commuting diagram
\[
\begin{array}{ccccccccc}
1&\longrightarrow&C_5&\longrightarrow&N_{12}
 &\longrightarrow&Q&\longrightarrow&1\\
 &&\Vert&&\downarrow\mathrm{fill}&&\Vert\\
1&\longrightarrow&C_5&\longrightarrow&N_{10}
 &\longrightarrow&Q&\longrightarrow&1 .
\end{array}                                                     \tag{F.14.8}
\]
Here $Q$ is the same subgroup of the mapping class
group of $X$: its elements preserve the covering
kernel and the two branch labels as a set, and
preserve the two free-puncture labels as a set.
The nonzero and zero monodromy values already
distinguish these two types.  The extra description
of a branch point as an original puncture on
$S_{12}$ therefore imposes no new permutation
condition.  On $S_{10}$ those same labels are retained
as branch decorations.  Descent of a filled
representative is the same homeomorphism of $X$,
which proves that both right vertical maps are
literally the identity after these identifications.
The kernel map is the named deck action itself.
The exact-sequence argument proves that the middle
filling map is an isomorphism.

Both finite groups in (F.14.8) are the same order-five
group.  The prime-order normalizer comparison and
full induced topology established in F.7.4--5 identify
their completed normalizers with their actual
closures in the ambient completions.  Thus completion
of the middle isomorphism is the actual injective
filling map on the completed normalizer.  In
particular it is injective on the source subgroup
from F.11 and on its index-two subgroup
$\widehat{\pi_1(Y)}$.
The ambient filling map is defined on the stabilizer
of the named pole pair; the normalizer preserves
that pair, and the source is pure, so this restriction
uses the stated ambient maps.

\emph{The commuting quotient square.}
The source still centralizes the actual involution
$\tau_t(z)=t/z$ after the poles are filled.
On $S_{10}$ this involution permutes the remaining
punctures in five pairs.  It has the same two
unmarked fixed points $s,-s$, labelled over $Y$.
Its decorated quotient has these two branch
points and five ordinary punctures.
The quotient map is still the literal function
$Q_s(z)$ in F.12.10.  The removed pole pair has
image $v_0=1$ and every remaining pair has the
same image $v_j$.

The completed involution-normalizer quotient again
has exactly its named finite deck kernel, by
F.7.4--5.  That deck involution permutes the ten
punctures nontrivially.  It has trivial intersection
with the completed pure subgroup, detected by
the finite permutation map.  Hence descent is
injective on the pure completed centralizer of
$\tau$ in the ten-puncture group.
The branch-labelled source lands in the pure
seven-puncture target.  Write $C^{\rm br}_{12}$
and $C^{\rm br}_{10}$ for the subgroups of the
two pure completed $\tau$-centralizers whose
descents also fix its two branch labels
individually.  Filling the pole pair preserves
this condition.
We have the actual commuting square
\[
\begin{array}{ccc}
 \widehat{\pi_1(Y)}
   &\longrightarrow&C^{\rm br}_{12}\\
 \big\downarrow\widehat\psi
   &&\big\downarrow\mathrm{fill}_{\{0,10\}}\\
 \widehat{P\Gamma}_{0,7}
   &\longleftarrow&C^{\rm br}_{10}.
\end{array}                                                     \tag{F.14.9}
\]
The upper arrow is the F.11--F.12 source inclusion;
the lower arrow is the just specified decorated
descent.  On that source the right arrow is
injective by (F.14.8), and the lower arrow is
injective on the pure subgroup.  Their composite
is precisely $(q_{v_0})_*\widehat\phi$, by the
pointwise formula for $Q_s$.  This proves
(F.14.2).

\emph{Both seven-point maps and their common coordinate.}
After forgetting $v_0$, the three labels $0,\infty,v_1$
give the unique normalization fixing $0,\infty$
and sending $v_1$ to one.  This gives the first
line of (F.14.3).
Substitution of F.12.4 followed by this same
normalization gives its second line.
Since $R(\phi(s_0))=\phi(s_0)$ with the exact
based representative already proved in F.12,
no additional target base path is required.
For clarity their common coordinate is
\[
 \frac{v_3(s)}{v_1(s)}
   =\left(
       \frac{(\zeta^2-s)(1+s)}
            {(\zeta^2+s)(1-s)}
     \right)^2 .
                                                               \tag{F.14.10}
\]
Thus forgetting all but its three normalization
labels and this coordinate gives equal composites
to $M_{0,4}$ as actual maps of spaces.
No injectivity of that four-point composite is used.

\emph{The branch-forgetting maps and based meridians.}
Forget the branch point $0$ instead, retaining
$\infty,v_0=1,v_1,\ldots,v_5$.
Normalize $v_0$ to zero and $v_1$ to one, keeping
infinity fixed.  Its four coordinates are
\[
 u_j(s)=\frac{v_{j+1}(s)-1}{v_1(s)-1}
       =\frac{\zeta^j(1+s)^2}{(\zeta^j+s)^2},
                \qquad 1\le j\le4.
                                                               \tag{F.14.11}
\]
They extend over $s=0$ and $s=\infty$, with respective
values $\zeta^{-j}$ and $\zeta^j$.  At either end
these four values are pairwise distinct and avoid
$0,1,\infty$.  Thus this is an extension into the
actual seven-point configuration space, not merely
a finite limit in its compactification.
The other branch-forgetting map has the same
property: the actual target coordinate inversion
$V\mapsto1/V$ exchanges the two branch labels and
satisfies $1/v_j(s)=v_j(-s)$.  Apply the same
normalization and extension with $-s$.
This proves the space-level factorization through
filling the two indicated source punctures and
hence (F.14.4) on full completions.

Here are explicit based choices.  Put
$\epsilon=1/10$ and use the polygonal whiskers
\[
 \begin{split}
 \omega_0&=(s_0,\,-2+2i,\,\epsilon+2i,\,\epsilon),\\
 \omega_\infty&=(s_0,\,-2+2i,\,10+2i,\,10).
 \end{split}                                                   \tag{F.14.12}
\]
From the exact expression
$s_0=\theta-\theta^3/2+i/2$ and $\theta>19/10$,
its real part is less than $-1$.
The first segments therefore lie to the left of
the unit circle; the horizontal segments have
imaginary part two.  The final segment of
$\omega_\infty$ lies outside that circle.
The final segment of $\omega_0$ has real part
$1/10$, different from zero and from every real
part of a point of $\pm\mu_5$.  Indeed their
nonzero real parts have absolute value at least
$(\sqrt5-1)/4>1/10$.
Both whiskers lie in $Y$.

Use the local circles
$s=\epsilon e^{2\pi i t}$ and $s=10e^{-2\pi i t}$,
respectively, and return along the reverse
whisker.  These give positive based source
meridians $m_0,m_\infty$.
The extended maps in (F.14.11) carry each local
circle to a loop filled by its actual disk
at the added source point; composing with the
image whiskers gives based null-homotopies.
The same holds for the other branch-forgetting
map.
In source abelianization $m_0$ is $e_0$ and
$m_\infty$ is $e_\infty=-e_0-\sum(e_j^++e_j^-)$.
A finite-puncture winding character takes each
of these meridians to $1$ or $-1$.
Its full continuous extension consequently
identifies the closure of each cyclic subgroup
with a full $\widehat\Z$ and detects every nonzero
profinite exponent.  Thus these are genuine
nontrivial full procyclic subgroups of both
branch-forgetting kernels.

\emph{The exact point-pushing defect.}
In the target labels of (F.14.1), the permutation
underlying $R$ exchanges $v_0$ with $v_2$.
Naturality of forgetting with this actual
labelled map gives
\[
                    R_*(K_0)=K_2.                            \tag{F.14.13}
\]
For $u\in\mathcal E_7$, the defining equality in
(F.14.5) says exactly that $d(u)\in K_0$.
Since $R_*^2=1$, direct multiplication gives
$R_*d(u)=d(u)^{-1}$.
Equation (F.14.13) then also puts $d(u)$ in
$K_2$.  The second identity in (F.14.6) follows
by expanding the definition on the product $uv$.
Both $K_i$ are normal, so the displayed conjugation
preserves $B$.  The third identity is the defining
full equality F.12.5, with both implications now
retained inside the explicitly specified
seven-point equalizer.

Finally the actual Birman point-pushing diagrams
identify $K_0$ with the completion of the
fundamental group of the sphere punctured at
the other seven labelled points, based at
$v_0$.  Forgetting $v_2$ restricts on this kernel
to the ordinary puncture-filling map
$\widehat F_6\to\widehat F_5$.
The point-pushing subgroup in the seven-point
target after forgetting $v_2$ is embedded.  The
actual completed Birman square is
\[
\begin{array}{ccccc}
 K_0=\widehat F_6&\longrightarrow&P_8
    &\longrightarrow&\widehat{P\Gamma}_{0,7}\\
 \downarrow\mathrm{fill}_{v_2}
    &&\downarrow\mathrm{forget}_{v_2}
    &&\downarrow\mathrm{forget}_{v_2}\\
 K'_0=\widehat F_5&\longrightarrow&
       \widehat{P\Gamma}_{0,7}
    &\longrightarrow&\widehat{P\Gamma}_{0,6}.
\end{array}                                                     \tag{F.14.14}
\]
The horizontal quotient maps forget $v_0$;
the lower middle group still retains $v_0$.
Consequently the kernel of that
restriction is exactly $K_0\cap K_2$.
The map of free groups kills the named peripheral
meridian about $v_2$ and keeps the other five
finite-puncture generators; the same finite
presentation and the universal property of
completion identify its full kernel with the
closed normal subgroup generated by that
meridian.

These full Birman maps can also be obtained
directly from normalized configuration spaces.
Forgetting one variable of an ordered configuration
in $\mathbf C\setminus\{0,1\}$ is a fibration
with the stated punctured-plane fibre, and has
a continuous section putting the forgotten
variable on the positive real axis beyond
$2+\sum|z_j|$.
The base and fibre are aspherical by induction
through these fibrations, giving the actual
discrete exact sequences.
The split group extension induces the full
profinite topology on its finitely generated
free kernel: use its characteristic finite-index
subgroups and the finite action on each resulting
finite quotient in the semidirect product.
Completion is then exact, and the forgetting
squares are the completions of the displayed
topological squares.
This justifies the kernel comparison on the
full completions used above.

The executed certificate
\path{code/part_a/common_reflection_forgetting_f14.py}
checks the two normalized composite formulas,
their common coordinate, the two endpoint
extensions, the exact source basepoint formula,
and the fourteen-by-twenty matrix of the
ordinary-point forgetting map on winding
characters.  The injectivity in (F.14.2)
uses the actual normalizer and covering maps,
not these homology checks.
\end{proof}
